\Subhead{4.11.1}{Линейный электрооптический и магнитооптический эффекты и их применение для управления светом}

Линейный электрооптический эффект меняет тензор оптической непроницаемости:
\begin{equation}
 \Delta(\varepsilon_r^{-1})_{ij}=r_{ijk}E_k.
\end{equation}
Он разрешён в соответствующих нецентросимметричных кристаллах. Для одной главной компоненты $\Delta n\simeq-n^3rE/2$; измеряемая разность фаз определяется \emph{разностью} изменений показателей двух собственных поляризаций. Пусть эффективный коэффициент определён как $\Delta n_{12}=n^3r_{\rm eff}E/2$. Тогда в поперечной схеме с зазором $d$ и длиной $L$
\begin{equation}
 \Gamma=\frac{2\pi L}{\lambda_0}\Delta n_{12}
 =\frac{\pi n^3r_{\rm eff}LV}{\lambda_0d},\qquad
 V_\pi=\frac{\lambda_0d}{n^3r_{\rm eff}L}.
\end{equation}
В продольной схеме $L=d$ геометрическая зависимость сокращается при таком определении коэффициента. Реальные тензорные ориентации могут давать иной $r_{\rm eff}$ и знак.

Между скрещёнными поляризаторами при ориентации входной поляризации под $45^\circ$ к осям $I_{\rm out}=I_{\rm in}\sin^2(\Gamma/2)$. Для линейной малосигнальной модуляции выбирают рабочую точку с ненулевым наклоном, обычно около $\Gamma=\pi/2$. Сам кристалл прежде всего меняет фазу/поляризацию, анализатор преобразует это в интенсивность. На высоких частотах важны ёмкость электродов и согласование скорости управляющей волны со светом.

Эффект Фарадея возникает из различия показателей двух круговых поляризаций при магнитном поле вдоль распространения. Угол $\theta_F=V_{\rm V}\int B_\parallel\dd l$, где $V_{\rm V}$ — постоянная Верде. При обратном проходе вращение в лабораторной системе складывается, в отличие от взаимного естественного оптического вращения. Фарадеевский ротатор с поляризаторами образует изолятор: обратная волна попадает на запрещённую поляризацию. Поглощение, эллиптичность и спектральная зависимость ограничивают полосу. Управление направлением пучка возможно через градиент создаваемого фазового набега.
